% =====================================================================
%  Thème 5 — Angle et identités remarquables — FACILE — Corrigé
% =====================================================================
\documentclass[11pt,a4paper]{article}
\usepackage[utf8]{inputenc}
\usepackage[T1]{fontenc}
\usepackage{lmodern}
\usepackage[french]{babel}
\usepackage{amsmath,amssymb,mathtools}
\usepackage{xcolor}
\usepackage{enumitem}
\usepackage{fancyhdr}
\usepackage[a4paper,margin=2.1cm,headheight=26pt,headsep=10pt]{geometry}

\definecolor{primary}{RGB}{222,70,80}
\definecolor{primarydark}{RGB}{150,30,42}
\definecolor{excol}{RGB}{0,135,90}

\pagestyle{fancy}\fancyhf{}
\fancyhead[L]{\small\textcolor{primarydark}{\textbf{Fiche 2} — Calcul vectoriel}}
\fancyhead[R]{\small\textcolor{primarydark}{\textsc{Thème 5 — Facile — Corrigé}}}
\fancyfoot[C]{\small\thepage}
\renewcommand{\headrulewidth}{0.8pt}

\newcommand{\vc}[1]{\overrightarrow{#1}}
\providecommand{\norm}[1]{\left\lVert #1\right\rVert}
\newcommand{\cor}[1]{\medskip\noindent{\color{primarydark}\bfseries\large Corrigé #1.}\par\smallskip}
\newcommand{\rep}{\textcolor{excol}{$\blacktriangleright$}\ }

\begin{document}
\begin{center}
{\LARGE\bfseries\color{primarydark}Angle et identités remarquables — Corrigé Facile}\\[2pt]
{\color{primary}\rule{0.55\linewidth}{1.1pt}}
\end{center}
\vspace{2mm}

\cor{1}
\begin{enumerate}[label=\textbf{\alph*)},leftmargin=2em,itemsep=2pt]
  \item $\vec u\cdot\vec v=1$, $\norm{\vec u}=1$, $\norm{\vec v}=\sqrt2$ :
        $\cos\theta=\dfrac{1}{\sqrt2}=\dfrac{\sqrt2}{2}$ (soit $45^\circ$).
  \item $\vec u\cdot\vec v=1$, $\norm{\vec u}=\sqrt{1+3}=2$, $\norm{\vec v}=1$ :
        $\cos\theta=\dfrac12$ (soit $60^\circ$).
  \item $\vec u\cdot\vec v=0$ : $\cos\theta=0$ (soit $90^\circ$, vecteurs perpendiculaires).
\end{enumerate}

\cor{2}
\begin{enumerate}[label=\textbf{\alph*)},leftmargin=2em,itemsep=2pt]
  \item $\cos\theta=\dfrac{\sqrt2}{2}\Rightarrow\theta=45^\circ$.
  \item $\cos\theta=\dfrac12\Rightarrow\theta=60^\circ$.
  \item $\cos\theta=\dfrac{\sqrt3}{2}\Rightarrow\theta=30^\circ$.
  \item $\cos\theta=0\Rightarrow\theta=90^\circ$.
\end{enumerate}
\rep Ces quatre valeurs sont à connaître par cœur.

\cor{3}
\begin{enumerate}[label=\textbf{\alph*)},leftmargin=2em,itemsep=2pt]
  \item $\vec u+\vec v\,(4;6)$, $\norm{\vec u+\vec v}=\sqrt{16+36}=\sqrt{52}=2\sqrt{13}$.
  \item $\vec u+\vec v\,(1;2)$, $\norm{\vec u+\vec v}=\sqrt{1+4}=\sqrt5$.
\end{enumerate}
\rep La méthode directe (additionner les composantes, puis prendre la norme) est la plus sûre.

\cor{4}
\begin{enumerate}[label=\textbf{\alph*)},leftmargin=2em,itemsep=2pt]
  \item $\norm{\vec u}^2=1^2+1^2=2$ ; $\norm{\vec v}^2=2^2+3^2=13$ ;
        $\vec u\cdot\vec v=1\times2+1\times3=5$.
  \item $\norm{\vec u+\vec v}^2=2+13+2\times5=25$. (On vérifie : $\vec u+\vec v\,(3;4)$,
        $3^2+4^2=25$.)
\end{enumerate}

\cor{5}
\begin{enumerate}[label=\textbf{\alph*)},leftmargin=2em,itemsep=2pt]
  \item $\norm{\vec u}^2=3^2+1^2=10$ ; $\norm{\vec v}^2=1^2+2^2=5$ ; d'où
        $\norm{\vec u}^2-\norm{\vec v}^2=10-5=5$.
  \item $(\vec u+\vec v)\cdot(\vec u-\vec v)=\norm{\vec u}^2-\norm{\vec v}^2=5$.
        Direct : $\vec u+\vec v\,(4;3)$, $\vec u-\vec v\,(2;-1)$, produit $=4\times2+3\times(-1)=8-3=5$.
        \checkmark
\end{enumerate}

\end{document}
